% @Author: Mouna REKIK
% @Date:   2024-10-11 / 12:30
%%%%%%%%%%%%%%%%%%%%%%%%%%%%%%%%%%%%%%%%%%%%%%%%%%%%%%%%%%%%%%%%%%%%%%%%%%%%%%

% \chapter*{REMERCIEMENT}
%\addstarredchapter{REMERCIEMENT}
%\addcontentsline{toc}{chapter}{REMERCIEMENT}
%\adjustmtc
% \justifying
% \sloppy
\newpage
\thispagestyle{MyStyle}
\vspace*{1cm}
\begin{center}{\fontfamily{ppl}\selectfont
\Huge{\textit{\textbf{Remerciements}}}}
\end{center} \vspace*{0.5cm}
\begin{center}
\itshape
~\par À l’issue de ce projet, je tiens à exprimer ma gratitude à toutes les personnes qui, par leurs conseils et leur soutien, ont contribué à sa réussite.\\
~\par Je remercie chaleureusement Madame Nouha BACCOUR, Maître Assistante à l’École Nationale d’Ingénieurs de Sfax (ENIS), pour son encadrement exemplaire, sa disponibilité, ses recommandations pertinentes et son accompagnement constant. Sa rigueur et ses orientations précises m’ont permis d’aborder ce travail avec méthode et de progresser efficacement.\\
~\par Je tiens également à adresser mes vifs remerciements à mes tuteurs industriels, Monsieur Omar ABID et Monsieur Omar BACCAR, pour l’environnement de travail favorable qu’ils m’ont offert, leur écoute, leurs conseils avisés, ainsi que leur patience tout au long de ce projet.\\
~\par Mes remerciements vont aussi à la société BACAB Consulting pour l’accueil bienveillant et l’opportunité unique de réaliser mon stage de fin d’études dans un cadre dynamique et formateur. Je salue l’ensemble de ses collaborateurs pour leur gentillesse, leur disponibilité et leur esprit collaboratif.\\
~\par Finalement, je tiens à témoigner ma profonde reconnaissance aux membres du jury pour l’honneur qu’ils me font en évaluant ce travail. Je souhaite que ce mémoire reflète l’engagement, la clarté et la rigueur que j’ai cherché à y insuffler.\\[2cm]
\end{center}
\begin{flushright}
\textbf{Omar CHABCHOUB}\\
Septembre 2025
\end{flushright}
